\documentclass[11pt,letterpaper]{article}
\usepackage[T1]{fontenc}
\usepackage[utf8]{inputenc}
\usepackage{lmodern}
\usepackage[margin=1in]{geometry}
\usepackage{booktabs}
\usepackage{longtable}
\usepackage{enumitem}
\usepackage{listings}
\usepackage{xcolor}
\usepackage[hidelinks]{hyperref}

\definecolor{backcolour}{rgb}{0.96,0.96,0.94}
\definecolor{cautionred}{rgb}{0.55,0,0}
\definecolor{okgreen}{rgb}{0,0.45,0}
\definecolor{unverifiedorange}{rgb}{0.75,0.35,0}

% See BASIC_CAN_COMMANDS.tex for why these options are set: everything in a
% listing must survive being selected from the PDF and pasted into a terminal.
\lstset{
    basicstyle=\ttfamily\small,
    backgroundcolor=\color{backcolour},
    numbers=none,
    frame=single,
    framesep=5pt,
    columns=fullflexible,
    keepspaces=true,
    upquote=true,
    showstringspaces=false,
    breaklines=false,
    xleftmargin=0pt,
    aboveskip=8pt,
    belowskip=8pt
}

\newenvironment{caution}[1]{%
    \par\addvspace{8pt}%
    \noindent{\color{cautionred}\rule{\linewidth}{0.8pt}}\par\vspace{2pt}
    \noindent\textbf{\color{cautionred}Caution: #1}\par\vspace{2pt}
}{%
    \par\vspace{2pt}\noindent{\color{cautionred}\rule{\linewidth}{0.4pt}}%
    \par\addvspace{8pt}%
}
\newcommand{\verified}{{\color{okgreen}\textbf{[VERIFIED]}}}
\newcommand{\unverified}{{\color{unverifiedorange}\textbf{[UNVERIFIED]}}}
\newcommand{\passif}[1]{\par\vspace{2pt}\noindent\textbf{\color{okgreen}Pass:} #1\par\vspace{4pt}}
\newcommand{\failif}[1]{\noindent\textbf{If it fails:} #1\par\vspace{6pt}}

\title{NEMA 23 / MKS SERVO57D --- Bring-Up SOP \\ \large Arctos Robotic Arm --- joint X}
\date{September 2026}

\begin{document}
\maketitle

\begin{caution}{This procedure has never been completed on real NEMA 23 hardware}
No NEMA 23 joint on this arm has ever communicated over CAN. The MKS SERVO57D fitted
to Joint X turned out to be the \textbf{RS485 variant}, which has no CAN transceiver,
so every step below that involves the CAN protocol is extrapolated from the
SERVO42D work rather than confirmed on a 57D.

Each step is tagged \verified{} or \unverified{}. \verified{} means confirmed on
this arm and expected to carry over directly, which is generally the host-side and
physical steps. \unverified{} means assumed by analogy with the 42D --- plausible,
but nobody has seen it work on a 57D. Confirm each of those as you go and update
this document with what you actually observe.
\end{caution}

Background, protocol details and troubleshooting live in
\texttt{Arctos\_CAN\_Interface.pdf}. Bare terminal commands live in
\texttt{BASIC\_CAN\_COMMANDS.pdf}.

\section*{The joint}

\begin{center}
\begin{tabular}{ll}
\toprule
\textbf{Joint} & X, base and waist rotation \\
\textbf{CAN ID} & \texttt{0x01} per config; confirm by scan \\
\textbf{Gear ratio} & 13.5:1 \\
\textbf{Motor} & NEMA 23, 1.8~N\textperiodcentered m, 2.8~A, 6.35~mm shaft \\
\textbf{Driver} & MKS SERVO57D \\
\bottomrule
\end{tabular}
\end{center}

\begin{caution}{A NEMA 23 is substantially more dangerous than a NEMA 17}
1.8~N\textperiodcentered m through 13.5:1 is roughly \textbf{24~N\textperiodcentered m
at the joint}, about 75 times the NEMA 17 joints' output torque, and more than enough
to injure a hand or destroy a printed part before you can react. Keep clear of the
mechanism, keep the emergency stop pasted in a second terminal, and keep a hand on
the power switch for every motion step.
\end{caution}

\section*{Step 0 --- Confirm the board is the CAN variant \verified}

\textbf{Do this first. It is the failure that has already cost this project weeks.}

\begin{enumerate}[leftmargin=*]
    \item Read the part number and silkscreen on the MKS SERVO57D board.
    \item Confirm it is the \textbf{CAN} variant, not RS485. The two look similar and
          both present a differential pair.
\end{enumerate}

\begin{caution}{RS485 boards will never work, at any bitrate or address}
RS485 carries MKS's own async serial protocol, not CAN framing. A CAN controller and
an RS485 transceiver cannot talk over the same wires. If this board is RS485 you must
either swap it for the CAN variant, which keeps this whole procedure and every script
valid, or treat the joint as a separate RS485 subsystem with its own adapter and
driver software, which is out of scope here.
\end{caution}

\passif{Board is confirmed to be the CAN variant.}
\failif{Stop. No amount of software work will make an RS485 board appear on the bus.}

\section*{Step 1 --- Pre-power electrical checks \verified}

\begin{enumerate}[leftmargin=*]
    \item Everything powered off, measure CAN\_H to CAN\_L. Expect about 60~$\Omega$.
    \item Confirm only CAN\_H, CAN\_L and GND run between the CANable and the
          controller.
    \item \textbf{Check the supply.} A NEMA 23 at 2.8~A draws far more than the NEMA
          17 joints. Confirm the supply is rated for it and that the wiring to the
          driver is sized appropriately.
    \item Confirm the joint is mechanically clear and that you know its travel limits.
\end{enumerate}

\passif{About 60~$\Omega$, correct wiring, adequate supply.}

\section*{Step 2 --- Configure the controller panel \unverified}

The 57D's menu layout has not been confirmed to match the 42D's. Set and
\textbf{save}:

\begin{center}
\begin{tabular}{ll}
\toprule
CAN ID & \texttt{0x01} \\
CAN Rate & \texttt{500} \\
CAN Response, shown as CANRSP & \texttt{ENABLE} \\
CAN CRC & \texttt{SUM-8} \\
Working current & 2.8~A or below, the NEMA 23's rating \\
\bottomrule
\end{tabular}
\end{center}

Power-cycle after saving.

\begin{caution}{Set the working current deliberately}
A 42D in this project was found storing 3000~mA against a 1.2~A motor. On a 2.8~A
NEMA 23 the consequences of an over-set current are correspondingly larger.
\end{caution}

\passif{Settings survive a power cycle.}
\failif{If the menu differs from the above, record what it actually offers and update
this document. That is exactly the unverified part.}

\section*{Step 3 --- Bring up the CAN interface \verified}

Identical to the NEMA 17 procedure. This is host-side and motor-independent.

\begin{lstlisting}
lsusb
ls -l /dev/serial/by-id/
sudo pkill slcand
sudo ip link delete can0 2>/dev/null
sudo slcand -o -c -s6 /dev/serial/by-id/usb-Openlight_Labs_CANable2* can0
sudo ip link set can0 txqueuelen 1000
sudo ip link set up can0
ip -br link show can0
\end{lstlisting}

\passif{The last command prints \texttt{can0} followed by \texttt{UP}.}

\section*{Step 4 --- Confirm the real CAN ID \unverified}

The scan itself is verified; what is unverified is whether a 57D answers the
\texttt{0x31} query the same way a 42D does.

\begin{lstlisting}
cd ~/arctos_ws/src/arctos_can_control/arctos_can_control
python3 can_id_scan.py --channel can0 --ids 1-16
\end{lstlisting}

Read-only. It never enables coils and never commands motion.

\passif{Exactly one responder, with a valid checksum.}
\failif{\textbf{Total silence with the board powered and the wiring confirmed is the
RS485 signature.} Go back to Step 0 and check the part number before doing anything
else. Note that the slcan adapter does not populate the kernel's bus-error counters,
so you cannot use ``zero errors'' as corroboration here.}

\section*{Step 5 --- Read-only health check \unverified}

\begin{lstlisting}
python3 joint_reliability_test.py --channel can0 --can-id 0x01 \
    --gear-ratio 13.5 --samples 200
\end{lstlisting}

\begin{center}
\begin{tabular}{lll}
\toprule
\textbf{Query} & \textbf{Opcode} & \textbf{Expected, by analogy} \\ \midrule
Encoder          & \texttt{0x31} & stable; zero jitter, zero drift \\
Velocity         & \texttt{0x32} & zero \\
Position error   & \texttt{0x39} & small \\
Enable state     & \texttt{0x3A} & \texttt{01} \\
Shaft protection & \texttt{0x3E} & \texttt{00} \\
\bottomrule
\end{tabular}
\end{center}

\begin{caution}{Confirm the encoder resolution on the 57D}
This procedure assumes 16384 counts per revolution and 3200 microstep pulses per
revolution, as on the 42D. \textbf{Verify this before trusting any angle
arithmetic}: command a known number of pulses and check that the encoder delta
matches. If the 57D differs, every conversion in this SOP changes.
\end{caution}

\passif{100\% response, zero checksum failures, zero jitter at rest.}

\section*{Step 6 --- Check the mechanism by hand \verified}

\begin{enumerate}[leftmargin=*]
    \item Disable the coils:
\begin{lstlisting}
cansend can0 001#F300F4
\end{lstlisting}
    \item Confirm the enable state now reads \texttt{00}:
\begin{lstlisting}
cansend can0 001#3A3B
\end{lstlisting}
    \item Turn the joint output by hand, gently, both directions, watching the
          encoder.
\end{enumerate}

\begin{caution}{Do not drop coils on a loaded base joint}
Joint X carries the entire arm above it. Removing holding torque may let the whole
assembly swing. Only do this with the arm supported or removed.
\end{caution}

\passif{The encoder tracks your hand both ways and the motion feels smooth.}
\failif{Gritty resistance means the gear mesh needs a plastic-safe grease. No encoder
change while the output turns means the motor is not coupled to it.}

\section*{Step 7 --- First commanded motion \unverified}

Paste the emergency stop into a second terminal, ready but unsent, before sending
anything:

\begin{lstlisting}
cansend can0 001#F7F8
\end{lstlisting}

Then one small move:

\begin{lstlisting}
python3 joint_unstick_move_test.py --can-id 0x01 --gear-ratio 13.5 \
    --step-degrees 0.25 --num-steps 1 --max-unstick 0 \
    --band-lo -1 --band-hi 1 --speed 25
\end{lstlisting}

\begin{caution}{Start smaller than you would on a NEMA 17}
At 13.5:1 the reduction is far lower than the NEMA 17 joints' 48:1 and 67.82:1, so
the same commanded joint angle turns the motor far less, and the joint moves much
faster for a given motor speed. A 0.25\textdegree{} joint step here is only about
3.4\textdegree{} of motor rotation. Verify direction and magnitude on this first
step before increasing anything.
\end{caution}

\passif{About 100\% of commanded, with \texttt{FD=[1, 2]}.}
\failif{Do not retry into a stall. Return to Step 6.}

\section*{Step 8 --- Stepped travel test \unverified}

\begin{lstlisting}
python3 joint_unstick_move_test.py --can-id 0x01 --gear-ratio 13.5 \
    --step-degrees 2.0 --num-steps 5 --max-unstick 0 \
    --band-lo -15 --band-hi 15 --speed 25
\end{lstlisting}

Then the same distance back:

\begin{lstlisting}
python3 joint_unstick_move_test.py --can-id 0x01 --gear-ratio 13.5 \
    --step-degrees -2.0 --num-steps 5 --max-unstick 0 \
    --band-lo -15 --band-hi 15 --speed 25
\end{lstlisting}

\passif{Every step 99--101\% with \texttt{FD 02}, and the round trip returns to
start.}

\section*{Step 9 --- Establish this motor's speed limit \unverified}

\textbf{The NEMA 17 figures do not transfer.} The 57D has different winding
inductance and a different torque-speed curve, so this motor's usable ceiling must
be measured rather than assumed. The measured table in
\texttt{Arctos\_CAN\_Interface.pdf} is 42D data; do not read Joint B's numbers
across to a 57D.

With the joint unloaded and clear:

\begin{lstlisting}
python3 joint_speed_limit_test.py --can-id 0x01 --gear-ratio 13.5 \
    --arc-lo -12 --arc-hi 12 --band-lo -15 --band-hi 15 \
    --speeds 50,100,150,200,250 --accel 8 --pause-s 15
\end{lstlisting}

The usable limit is the highest commanded speed where actual rpm still tracks it to
within about 1\%.

\begin{caution}{Keep the duty cycle down}
Running a long ladder of trials back to back overheats the motor. Fifteen seconds
between trials is reasonable; one second is not. Stop as soon as you have the
answer, and put a hand on the motor part way through --- there is no temperature
opcode, so nothing in software will warn you.
\end{caution}

Record the result here once measured:

\begin{center}
\begin{tabular}{ll}
\toprule
Tracks commanded within 1\% up to & \underline{\hspace{3cm}} rpm \\
Roll-off begins                   & \underline{\hspace{3cm}} rpm \\
Saturation ceiling                & \underline{\hspace{3cm}} rpm \\
\bottomrule
\end{tabular}
\end{center}

\section*{Step 10 --- Acceptance}

\begin{center}
\begin{tabular}{ll}
\toprule
\textbf{Check} & \textbf{Criterion} \\ \midrule
Board variant     & confirmed CAN, not RS485 \\
CAN ID            & confirmed by scan \\
Comms             & 100\% response, zero checksum failures \\
Encoder at rest   & zero jitter, zero drift \\
Encoder scale     & confirmed for the 57D, not assumed \\
Hand check        & smooth, tracks both directions \\
Per-step accuracy & 99--101\% of commanded \\
Status frames     & \texttt{FD 02} on every step \\
Round trip        & returns to start \\
Speed limit       & measured and recorded above \\
\bottomrule
\end{tabular}
\end{center}

\section*{Step 11 --- Shutdown \verified}

Stop any running script with Ctrl+C first.

\begin{lstlisting}
cansend can0 001#F300F4
sudo ip link set can0 down
sudo pkill slcand
\end{lstlisting}

Then unplug the adapter. Skip the first line if the arm above this joint is not
supported and needs it to keep holding.

\section*{Update this document}

Every step tagged \unverified{} is a claim nobody has tested on a 57D. When you get
CAN-variant hardware onto the bus, work through them in order and replace the tag
with what you actually saw, including anything that turned out to differ from the
42D.

\end{document}
